\chapter*{Abstract}
\addcontentsline{toc}{chapter}{Abstract}

The textile industry relies heavily on proprietary Computer-Aided Design (CAD) software, such as NedGraphics, Pointcarré, and ArahWeave, to transition from creative concepts to loom-ready manufacturing data. However, these traditional systems carry prohibitive annual licensing costs ranging from \$3,000 to \$8,000 for NedGraphics and Pointcarré, with ArahWeave Personal Edition priced at approximately €2,500 per year (~\$2,750). they feature incredibly complex, grid-centric user interfaces that require specialist training, creating a massive barrier to entry. This digital divide is particularly evident in the Indian handloom sector, enveloping weaving clusters in Varanasi, Surat, and Bhiwandi, where independent artisans cannot afford to digitize their workflow. Despite the rapid advancement of Artificial Intelligence, there currently exists no AI-native CAD tool in the textile domain that accepts intuitive, natural language design intent and translates it into functional mathematical weave structures.

Dobby Studio is an AI-powered, completely browser-based fabric design platform built utilizing a modern React 18 and Vite frontend, integrated seamlessly with the Groq Large Language Model API. The system utilizes `llama-3.3-70b-versatile` to enable instantaneous natural language to fabric sett generation, completely bypassing traditional pixel-drawing workflows. The application performs real-time 2D fabric rendering utilizing an advanced optical color-mixing algorithm that precisely mimics the physical 70/30 warp and weft blending. To aid designers, it provides a 3D cloth drape simulation via Three.js and generates production-grade export artifacts including JSON, PNG, printable specification sheets, and industry-standard WIF 1.1 formatted files. Crucially, the system implements mathematically accurate reflective sett repeat logic and supports dynamic Ends Per Inch (EPI) and Picks Per Inch (PPI) physical size calculations.

The implemented system successfully eliminates the immense financial and technical barriers previously hindering freelance textile designers and SME manufacturers. System evaluations demonstrate highly successful outcomes, achieving 91\% accuracy in AI-driven text-to-sett generation and guaranteed 100\% validity in its WIF export engine for loom compatibility. Deployed globally on the Vercel architecture, Dobby Studio maintains sub-second response times with modal latency of 800–1200 milliseconds while functioning entirely free of charge. By providing an interface accessible via any web browser, the platform acts as a modern, democratized alternative to commercial CAD software, directly targeting India's 35 lakh handloom weavers and small weaving mills that were previously excluded from the digital design revolution.
